\section{Tableaux for $\Hxpath$} \label{sec:tableaux}

In this section we present the notion of tableaux for $\Hxpath$, comment on the view of these tableaux as a calculus for satisfiability, and prove the soundness and completeness of this calculus.
Our ideas and results build on the terminating tableau calculus for Hybrid Modal Logic in~\cite{BolanderB07}.
In brief, we incorporate to such a calculus rules for handling data (in)equalities $\tup{\alpha \cmpr \beta}$ and $[\alpha \cmpr \beta]$.
%The main challenge is how to handle node expressions of the form $\tup{\alpha \cmpr \beta}$ and $\lnot\tup{\alpha \cmpr \beta}$.
We consider our presentation of tableaux for $\Hxpath$ simpler and more elegant than the one presented in~\cite{ArecesFS17}.
First, the internalization of rules in our calculus simplifies their manipulation.
Second, our way of handling the left-hand-side of data (in)equalities together with a commutation rule significantly reduces the number of rules.
For the remainder of the article, we assume that $\noms$ is the set of natural numbers.

\subsection{Tableau Calculus.}
\label{subsec:rules}

We begin with the definition of a tableau for a node expression.
We assume familiarity with the usual terminology on labelled trees.
In particular, the  notion of a root, a node, a successor of a node, a branch, a leaf, and the label of a node.
Moreover, we assume some familiarity with tableaux for Hybrid Logic (see~\cite{BolanderB07}).

\begin{defi}\label{def:tableau}
    Let $\varphi$ be a node expression and $i$ a new nominal.
    A \emph{tableau} for $\varphi$ is a labelled tree constructed from an initial single node tree labelled with $i{:}\varphi$ using the branch \emph{extension} rules in \Cref{rules:basic,rules:nominals,rules:internalization,rules:path}. Crucially, rules can be applied only once to its premisses. 
    In this definition, a \emph{new} nominal is the smallest nominal that is larger than all other nominals appearing in the tableau (thus far constructed).
    In turn, a \emph{root} nominal is one appearing in $\varphi$.
    Finally, a node expression $i{:}\tup{\pathh}j$ is an \emph{accessibility constraint} iff it appears in the conclusion of the rules ($\Diamond$) and (child).
\end{defi}

%BASIC 

\begin{figure}[p]% BASIC
    \centering
    \small
    \begin{tabular}{ccccc}
        \toprule
        \multicolumn{5}{c}{\sc Basic}
        \tabularnewline
        \midrule
        {\infer[(\lnot\lnot)]
            {\deduce[]
                {\phantom{|}}
                {\phantom{|}i{:}\varphi}}
            {\deduce[]
            {\phantom{|}{i{:}\lnot \lnot \varphi}}
            {}
            }      
        }
        &    
        {\infer[(\land)]
            {\deduce[]
                {\phantom{|}i{:}\varphi}
                {\phantom{|}i{:}\psi}}
            {\deduce[]
            {\phantom{|}i{:}(\varphi \land \psi)}
            {}
            }
        }
        &
        {\infer[(\lnot \land)]
            {\deduce[]
            {\phantom{|}}
            {\phantom{|}{i{:}\lnot\varphi} \quad {i{:}\lnot\psi}}
            }
            {\deduce[]
            {\phantom{|}{i{:}\lnot(\varphi \land \psi)}}
            {}
            }
        }
        &
        {\infer[(\Diamond^{1})]
            {\deduce[]
            {\phantom{|}{j{:}\varphi}}
            {\phantom{|}{i{:} \tup{\child} j}}
            }
            {\deduce[]
            {\phantom{|}{i{:} \tup{\child} \varphi}}    
            {\phantom{|}}
            }
        }
        &
        {\infer[(\lnot \Diamond^2)]
            {\deduce[]
            {\phantom{|}}
            {\phantom{|}{j{:} \lnot \varphi}}
            }
            {\deduce[]
            {\phantom{|} {i{:} \tup{\child} j}}    
            {\phantom{|}{i{:} \lnot \tup{\child} \varphi}}
            }
        }
        \tabularnewline
        \bottomrule
        \multicolumn{5}{l}{\footnotesize $^1$ $i{:}\tup{\child}\varphi$ is not an accessibility constraint. $j$ is a new nominal.}
        % \tabularnewline
        % \multicolumn{5}{l}{\footnotesize $^2$ $i{:}\tup{\child}\varphi$ is not an accessibility constraint.}
        \tabularnewline
        \multicolumn{5}{l}{\footnotesize $^2$ $i{:}\tup{\child}j$ is an accessibility constraint.}
    \end{tabular}
    \caption{Basic Rules}
    \label{rules:basic}
\end{figure}


\begin{figure}[p]% SATISFIABILITY AND NOMINALS
    \centering
    \small
    \begin{tabular}{cccccc}
        \toprule
        \multicolumn{6}{c}{\sc Satisfiability and Nominals}
        \tabularnewline
        \midrule
        {\infer[(\mathrm{nom})]
            {\phantom{|}{j{:}\varphi}}
            {\phantom{|}i{:} (j{:} \varphi)}
        }
        &
        {\infer[(\lnot \mathrm{nom})]
            {\phantom{|}{j{:} \lnot \varphi}}
            {\phantom{|}{i{:} \lnot (j{:} \varphi)}}
        }
        &         
        {\infer[(\mathrm{copy}^1)]
            {\phantom{|} j{:} \varphi }
            {\deduce[]
                {\phantom{|}i{:}j}
                {\phantom{|}i{:}\varphi}
            }
        }
        &
        {\infer[(\mathrm{ref}^2)]
            {\phantom{|} i{:}i}
            {\deduce[]
            {\phantom{|}}           
            {\deduce[]
            {\phantom{|}}
            {\phantom{|}} 
            }
            }
        }                          
        &
        {\infer[(\mathrm{sym})]
            {\phantom{|} i{:}j }
            {\deduce[]
            {\phantom{|}j{:}i}
            {\phantom{|}} 
            }
        }
        &
        {\infer[(\mathrm{trans})]
            {\phantom{|} j{:}k }
            {\deduce[]
            {\phantom{|}j{:}l}
            {\deduce[]
            {\phantom{|} i{:}l} 
            {\phantom{|} i{:}k} 
            }
            }                        
        }
        \tabularnewline
        \bottomrule
        \multicolumn{6}{l}{\footnotesize $^1$ $\varphi$ is not a nominal, $j$ is a root nominal, $j<i$.}
        \tabularnewline
        \multicolumn{6}{l}{\footnotesize $^2$ $i$ is a nominal in the branch.}
        % \tabularnewline
        % \multicolumn{6}{l}{\footnotesize }
    \end{tabular} %}
    \caption{Rules for Nominals}
    \label{rules:nominals}
\end{figure}

\begin{figure}[p]% INTERNALIZATION
    \centering
    \small
    % \resizebox*{.85\textwidth}{!}
    % {
    \begin{tabular}{cccc}
        \toprule
        \multicolumn{4}{c}{\sc Internalization into Data (In)Equalities}
        \tabularnewline
        \midrule
        {\infer[(\mathrm{int}_1)]
            {{\phantom{|}}\langle \goto{i}\alpha \cmpr \goto{i}\beta \rangle}
            {i{:} \langle \alpha \cmpr \beta \rangle}
        }
        &
        {\infer[(\mathrm{int}_2)]
            {\phantom{|} \lnot \langle \goto{i}\alpha \cmpr \goto{i}\beta \rangle }
            {\phantom{|}i{:} \lnot \langle \alpha \cmpr \beta \rangle}
        }
        &
        {\infer[({\lnot}{\cmpr})]
            {\phantom{|}\langle \goto{i} \cmpd \goto{j} \rangle}
            {\phantom{|} \lnot \langle \goto{i} \cmpr \goto{j} \rangle} 
        }
        % &
        % {\infer[({\lnot}{=})]
        %     {\phantom{|}\langle \goto{n} \neq \goto{m} \rangle}
        %     {\phantom{|} \lnot \langle \goto{n} = \goto{m} \rangle}
        % }
        &
        {\infer[(\mathrm{dRef})]
            {\phantom{|} \langle \goto{i} =_\cmp \goto{j} \rangle }
            {\phantom{|} i{:}j}
        }
        % \tabularnewline[10pt]
        % \multicolumn{5}{c}%
        % {\infer[(\mathrm{dTrans})]
        %     {\langle \goto{m} = \goto{n}  \rangle}
        %     {\deduce[]
        %     { \langle k{:} = \goto{n}  \rangle }
        %     {\langle \goto{m} = k{:}  \rangle}}
        % }
        \tabularnewline
        \bottomrule
    \end{tabular} %}
    \caption{Rules for Internalizing Nominals and Negations into Data (In)Equalities}
    \label{rules:internalization}
\end{figure}

\begin{figure}[p]%PATH
    \centering
    \small
    \begin{tabular}{c@{\quad}c@{\quad}c}
        \toprule
        \multicolumn{3}{c}{\sc Path Expressions in Data (In)Equalities}
        \tabularnewline
        \midrule
        {\infer[(\mathrm{child}^1)]
            {\deduce[]      
            {\phantom{|}\langle \goto{j} \alpha \cmpr \beta \rangle } 
            {\phantom{|} i{:} \langle \child \rangle j  }     
            }    
            {\deduce[]
            {\phantom{|}\langle \goto{i} \child \alpha \cmpr \beta \rangle}
            {\phantom{|}} 
            }
        }
        &
        {\infer[(\lnot \mathrm{child})]
            {\deduce[]      
            {\phantom{|}} 
            {\phantom{|}  \lnot \langle \goto{j} \alpha \cmpr \beta \rangle  }     
            }    
            {\deduce[]
            {\phantom{|}i{:} \langle \child \rangle j }
            {\phantom{|}\lnot \langle \goto{i} \child \alpha \cmpr \beta \rangle} 
            }
        }
        &
        {\infer[(\mathrm{com}_1)]
            {\deduce[]
            {\phantom{|}}
            {\langle \beta \cmpr \goto{i} \rangle}}
            {\langle \goto{i} \cmpr \beta \rangle } 
        }
        \tabularnewline
        {\infer[(\mathrm{test})]
            {\deduce[]
            {\phantom{|}  \langle \goto{i} \alpha \cmpr \beta \rangle}
            {\phantom{|} i{:}\varphi }      
            }    
            {\deduce[]
            {\phantom{|} \langle \goto{i} \varphi? \alpha \cmpr \beta \rangle} 
            {\phantom{|} } 
            }  
        }
        &
        {\infer[(\lnot \mathrm{test} )]
            {\deduce[]
            {\phantom{|} }
            {\phantom{|}  i{:} \lnot \varphi \quad \lnot \langle \goto{i} \alpha \cmpr \beta \rangle }      
            }    
            {\deduce[]
            {\phantom{|} \lnot \langle \goto{i} \varphi? \alpha \cmpr \beta \rangle} 
            {\phantom{|} } 
            }  
        }
        &
        {\infer[(\mathrm{com}_2)]
            {\deduce[]
            {\phantom{|}}
            {\lnot\langle \beta \cmpr \goto{i} \rangle}}
            {\lnot\langle \goto{i} \cmpr \beta \rangle } 
        }
        % &
        % {\infer[(\mathrm{com}_2)]
        %     {\phantom{|}\lnot\langle \beta \cmpr \goto{i} \rangle}
        %     {\deduce[]
        %     {\phantom{|}\lnot\langle \goto{i} \cmpr \beta \rangle}
        %     {\phantom{|}}
        %     }
        % }
        \tabularnewline
        {\infer[(@)]
            {\phantom{|}  \langle \goto{j} \alpha  \cmpr \beta \rangle }
            {\deduce[]
            {\phantom{|} \langle \goto{i} \goto{j} \alpha \cmpr \beta \rangle} 
            {\phantom{|}} 
            }  
        }
        &
        {\infer[(\lnot @ )]
            {\phantom{|} \lnot \langle \goto{j} \alpha  \cmpr \beta \rangle }
            {\deduce[]
            {\phantom{|} \lnot \langle \goto{i} \goto{j} \alpha \cmpr \beta \rangle} 
            {\phantom{|}} 
            }  
        }
        &
        {\infer[(\mathrm{dTrans})]
            {\langle \goto{i} =_\cmp \goto{j}  \rangle}
            {\deduce[]
            { \langle \goto{k} =_\cmp \goto{j}  \rangle }
            {\langle \goto{i} =_\cmp \goto{k}  \rangle}}
        }
        \tabularnewline
        {\infer[(\cup)]
            {\deduce[]
            {\phantom{|}  }
            {\phantom{|} \langle \goto{i} \alpha\gamma \cmpr \eta \rangle \quad \langle \goto{i} \beta\gamma \cmpr \eta \rangle }      
            }    
            {\deduce[]
            {\langle \goto{i} (\alpha \cup \beta) \gamma \cmpr \eta \rangle} 
            {\phantom{|} } 
            }  
        }
        &
        {\infer[(\lnot \cup )]
            {\deduce[]
            {\phantom{|} \lnot \langle \goto{i} \beta\gamma \cmpr \eta \rangle }      
            {\phantom{|}  \lnot \langle \goto{i} \alpha\gamma \cmpr \eta \rangle }
            }    
            {\deduce[]
            {\phantom{|} \lnot \langle \goto{i} (\alpha \cup \beta) \gamma \cmpr \eta \rangle} 
            {\phantom{|}} 
            }  
        }
        \tabularnewline
        % &
        % {\infer[(\mathrm{com}_1)]
        %     {\langle \beta \cmpr \goto{n} \rangle}  
        %     {\langle \goto{n} \cmpr \beta \rangle } 
        %     }
        % &
        % {\infer[( \mathrm{com}_2 )]
        % {\lnot \langle \beta \cmpr \goto{n} \rangle}
        % {\lnot\langle \goto{n} \cmpr \beta \rangle } 
        % }
        % \tabularnewline
        \bottomrule
        % \tabularnewline
        \multicolumn{2}{l}{\footnotesize $^1$ $j$ is a new nominal.}
    \end{tabular} %}
    \caption{Rules for Handling Path Expressions in Data (In)Equalities}
    \label{rules:path}
\end{figure}


Let us briefly discuss how tableaux rules are meant to be understood. 
% in general and in particular.
In general, all rules are taken to be read as: if the premisses of the rule are in the branch, then, extend the branch by adding to its leaf node successors labelled with the conclusions of the rule. Each rule is applied only once to a given set of premisses present in a branch.
\Cref{rules:basic} presents the usual rules for Boolean and modal operators of Basic Modal Logic, with the distinguishing characteristic of accessibility constraints being internalized in the language using nominals.
\Cref{rules:nominals} presents standard rules for handling nominals, similar to those found in~\cite{BolanderB07,Brauner2011}.
The only difference is the rule (copy) whose side conditions are in place to guarantee termination of the calculus. 
%This is made clear in \Cref{sec:termination}.
\Cref{rules:internalization} introduces rules for data comparison, they internalize nominals into data (in)equalities.
More precisely, notice how the rule ($\mathrm{int}_1$) takes a satisfiability statement $i{:}\tup{\alpha \cmpr \beta}$, and internalizes $i{:}$ as an $@_i$-prefix in $\alpha$ and $\beta$ resulting in $\tup{@_i\alpha \cmpr @_i\beta}$. In doing this, we anchor the point of evaluation of the node expression at the beginning of each path in the data comparison.
This enables us to work out paths individually using the rules in~\Cref{rules:path}. The rule ($\mathrm{int}_2$) does the same job in presence of a negation. In turn, ($\neg{\cmpr}$) internalizes negated equalities (inequalities) into inequalities (equalities). 
%This can be viewed as a way of ``pushing negations inside'' data comparisons.
Finally, (dRef) tells us that nodes with the same name must have the same data values.
%This again may be viewed as a way of internalizing satisfiability statements into data equalities.
\Cref{rules:path} depicts the last set of rules in our tableaux.
These rules tell us how to deal with path expressions in data comparisons. 
%The rules (child) and ($\lnot$child) can be understood in analogy with the rules ($\Diamond$) and (${\lnot}{\Diamond}$).
Essentially, (child) mimics what is done in ($\Diamond$) but in the context of a data comparison, whereas ($\lnot$child) is analogous to (${\lnot}{\Diamond}$).
%The rules (test) and ($\lnot$test) can be understood in analogy with the rules ($\land$) and (${\lnot}{\land}$).
%Simply put, 
The rule (test) eliminates $\varphi$ from the path expression and evaluates it as a standalone node expression, while the case for ($\lnot$test) is dual. 
%It may be worth noting the correspondence between (test) and ($\lnot$test) and the properties in \Cref{prop:equivalences}. 
The rules ($@$) and (${\lnot}{@}$) can be understood in analogy with (nom) and (${\lnot}$nom).
The rationale behind ($\cup$) and (${\lnot}{\cup}$) is immediately understood from \Cref{prop:equivalences}. 
A novelty of our calculus is that rules only decompose, step-by-step, path expressions on the left side of a data comparison.
The treatment of the right side takes over after applying ($\mathrm{com}_1$) or  ($\mathrm{com}_2$). 
Finally, (dTrans) corresponds to transitivity of equalities.


%\resizebox{\textwidth}{!}
\begin{exa}\label{ex:usetableaux}
%\Cref{fig:atableau} depicts 
\Cref{fig:atableau} depicts a tableau for $\tup{\child}\tup{@_2\astep[b]2? =_\cmp \astep[b](q\land 3)?}$, starting from an initial node labelled by $4{:}\tup{\child}\tup{@_2\astep[b]2? =_\cmp \astep[b](q\land 3)?}$.
The annotations on the right indicate the rules we applied and the nodes we applied them to.
The tableau has a single branch, which becomes saturated if we apply the rules (ref) and (dRef) to all the nominals it contains --e.g., if we apply (ref) to $2$, we obtain a node 18 with label $2{:}2$, and if we apply (dRef) to the new node, we obtain a node 19 with label $\tup{@_2=_\cmp @_2}$; and so on for all the nominals in the branch.
\begin{figure}[t]
    \begin{center}
    \begin{prooftree}{}%
    [
        4{:}\tup{\child}\tup{@_2\astep[b]2? =_\cmp \astep[b](q\land 3)?},
        just={root}
        [
            4{:}\tup{\child}5, 
            just={($\Diamond$):!u}
            [
                5{:}\tup{ @_2\astep[b]2? =_\cmp \astep[b](q \land 3)?},
                just={($\Diamond$):!uu}
                [
                    \tup{@_5@_2\astep[b]2? =_\cmp @_5 \astep[b](q \land 3)?},
                    just={(int$_1$):!u}
                    [
                        \tup{@_2\astep[b]2? =_\cmp @_5 \astep[b](q \land 3)?},
                        just={(${@}$):!u}
                        [
                            2{:}\tup{\astep[b]} 6,
                            just={(child):!u}   
                            [   
                                \tup{@_62? =_\cmp @_5 \astep[b](q \land 3)?},
                                  just={(child):!u}
                                [
                                    6{:}2, 
                                    just={(test):!u}
                                    [
                                        2{:}6,
                                        just={(sym):!u}
                                        [
                                            \tup{@_6 =_\cmp @_5 \astep[b](q \land 3)?} ,
                                            just={(${\text{test}}$):!uuu}  
                                            [
                                                \tup{ @_5 \astep[b](q \land 3)? =_\cmp @_6} ,
                                                just={(com$_1$):!u} 
                                                [
                                                    5{:}\tup{\astep[b]}7, 
                                                    just={(child):!u}
                                                    [
                                                        \tup{ @_7 (q \land 3)? =_\cmp @_6},
                                                        just={(child):!uu}
                                                        [
                                                            7{:}(q \land 3), 
                                                            just={(test):!u}
                                                            [
                                                                \tup{ @_7 =_\cmp @_6},
                                                                just={(test):!uu}
                                                                [
                                                                    \tup{ @_6 =_\cmp @_7},
                                                                    just={(com$_1$):!u}
                                                                    [
                                                                        7{:}q, 
                                                                        just={($\land$):!uuu}
                                                                        [
                                                                            7{:}3, 
                                                                            just={($\land$):!uuuu}
                                                                            [
                                                                                3{:}7, 
                                                                                just={(sym):!u}
                                                                                [
                                                                                    3{:}q, 
                                                                                    just={(copy):!uuu,!uu}
                                                                                ]
                                                                            ]
                                                                        ]
                                                                    ]
                                                                ]
                                                            ]
                                                        ]
                                                    ]
                                                ]
                                            ]
                                        ]
                                    ]
                                ]
                            ]
                        ]
                    ]
                ]
            ]
        ]
    ]
    \end{prooftree}
\end{center}
   \caption{A tableau for $\tup{\child}\tup{@_2\astep[b]2? =_\cmp \astep[b](q\land 3)?}$}\label{fig:atableau}
\end{figure}
\end{exa}

We consider our tableaux to be simpler and more elegant than the proposal  in~\cite{ArecesFS17}. Treating only the left side of data comparisons significantly reduces  the number of rules. It is also clear that our tableaux are \emph{internalized}; they have (only) node expressions as labels. Internalized tableaux have important advantages. In particular, it will enable us to define extended calculi to handle many different frame conditions, and obtain completeness automatically as it is done in~\cite{Blackburn00,Blackburn2002BeyondPA} (see~\Cref{sec:pure-ext}). 
Our tableaux are also \emph{analytic} in the sense of~\Cref{def:analytic} below (see, e.g.,~\cite{Smullyan1968} for background on analytic tableaux). %; the formulation of which is done with the aid of \Cref{def:subformulas,def:quasi:component} detailing what we mean by the ``parts'' making up a node expression in a tableaux. %In order to define the notion of analytic, we need first to provide some definitions.
This property is helpful to provide a termination argument (see \Cref{sec:termination}).

\begin{defi} \label{def:subformulas}
    The \emph{subcomponents} of a node expression $\varphi$ are defined as
        \begin{longtable}{r@{\ }c@{\ }l}
            $\sub(p)$
            & $=$ & $\set{p}$ \\
            $\sub(i)$
            & $=$ & $\set{i}$ \\
            $\sub(\lnot \varphi)$
            & $=$ & $\set{\lnot \varphi} \cup \sub(\varphi)$ \\
            $\sub(\varphi \land \psi)$
            & $=$ & $\set{\varphi \land \psi} \cup \sub(\varphi) \cup \sub({\psi})$  \\            
            $\sub(i{:}\varphi)$
            & $=$ & $\set{i,\, i{:}\varphi} \cup \sub(\varphi)$ \\
            $\sub(\tup{\pathh} \varphi)$ 
            & $=$ & $\set{ \tup{\pathh} \varphi} \cup \sub(\varphi)$ \\  
            $\sub(\tup{ \patha \cmpr \patha[b] })$
            & $=$ & $\set{\tup{ \patha \cmpr \patha[b]}} \cup \sub'(\patha) \cup \sub'(\patha[b])$ \\
            $\sub(\tup{\patha \cmpr \patha[b]\beta})$
            & $=$ & $\set{\tup{\patha \cmpr \patha[b]\beta}} \cup \sub(\tup{\patha[b]\beta \cmpr \patha})$ \\
            $\sub(\tup{\patha\alpha \cmpr \beta})$
            & $=$ & $\set{\tup{\patha\alpha \cmpr \beta}} \cup \sub(\tup{\alpha \cmpr \beta}) \cup \sub'(\patha)$ \\
            $\sub(\tup{ (\alpha \cup \beta) \gamma \cmpr \xi })$
            & $=$ & $\set{\tup{ (\alpha \cup \beta) \gamma \cmpr \xi }} \cup \sub{(\tup{ \alpha \gamma \cmpr \xi })} \cup \sub{(\tup{ \beta \gamma \cmpr \xi })}$.
        \end{longtable}
    % where $\sub'$ is such that:
    % \begin{align*}
    %     \sub'(\goto{i}) & = \set{i} & \sub'(\child) & = \emptyset & \sub'(\varphi?) & = \sub(\varphi).
    % \end{align*}
    where $\patha \in \{\goto{i}, \pathh, \varphi?\}$. %; and $\pathaas$ indicates a path expression beginning with $\patha$ (with $\alpha$ indicating any well-formed path expression).
    $\sub'(\goto{i}) = \set{i}$, $\sub'(\child) = \emptyset$, and $\sub'(\varphi?) = \sub(\varphi)$.
\end{defi}

Subcomponents are the basic parts tableau rules operate on as they break down complex expressions.
During this process, the rules often prefix them with nominals, associating parts of expressions with specific points in a model. The definition of quasicomponents below formalizes which combinations of nominals and subcomponents can arise during this process.


\begin{defi} \label{def:quasi:component}
        The set of quasi subcomponents of a node expression
        $i{:}\varphi$ is defined as
    \begin{align*}
        \qsub(i{:}\varphi) = \ & \set{j{:}\psi, j{:}\lnot\psi \mid \psi\in\sub(i{:}\psi), j \in \noms}   \\
            &  \cup \{\tup{\goto{j}\alpha \cmpr \goto{k}\beta}, 
            \tup{\goto{j} \cmpr \goto{k}\beta}, 
            \tup{\goto{j}\alpha \cmpr \goto{k}}, 
            \tup{\goto{j} \cmpr \goto{k}} \\ 
            & \quad \ \mid \tup{\alpha \cmpr \beta} \in \sub(i{:}\varphi), j \in \noms\} \\
        & \cup \{\lnot \tup{\goto{j}\alpha \cmpr \goto{k}\beta},
        \lnot \tup{\goto{j} \cmpr \goto{k}\beta},
        \lnot \tup{\goto{j}\alpha \cmpr \goto{k}},
        \lnot \tup{\goto{j} \cmpr \goto{k}},
        \tup{\goto{j} \cmpd \goto{k}}  \\ 
        & 
        \quad \ \mid  \lnot \tup{\alpha \cmpr \beta} \in \sub(i{:}\varphi), j \in \noms\}.
    \end{align*}
\end{defi}

The cases where the quasi subcomponent are $j{:}\psi$, or $j{:}\lnot\psi$ 
arise naturally from the basic tableau rules. %, including those for satisfiability and nominals.
The other cases, correspond to applications of commutativity rules and for handling data equalities.
For instance, $\tup{\goto{j}\alpha \cmpr \goto{k}\beta}$ and $\lnot\tup{\goto{j}\alpha \cmpr \goto{k}\beta}$ may arise as the result of applying one of the internalization rules in \Cref{rules:internalization}.
The case $\tup{\goto{j} \cmpr \goto{k}\beta}$ corresponds to a situation in which the left-hand side of a data (in)equality has been fully decomposed. 
Similarly, $\tup{\goto{j} \cmpr \goto{k}}$ corresponds to a situation in which both sides of a data equality have been fully decomposed, and
 $\tup{\goto{j} \cmpd \goto{k}}$ may arise from applying (${\lnot}{\cmpr}$).


\begin{defi} \label{def:analytic}
    A tableau for a node expression $\varphi$ is \emph{analytic} iff the label of any node in the tableau is one of the following:
    \begin{enumerate}
        \item a quasi subcomponent of the root,
        \item an accessibility constraint,
        \item an equality constraint $\tup{@_i =_\cmp @_j}$ resulting from the rules (dRef) and (dTrans), or
        \item a satisfiability constraint $i{:}j$ resulting from the rules (ref), (sym), or (trans).
    \end{enumerate}
\end{defi}

\begin{lem} \label{prop:analytic}
    The tableaux in \Cref{def:tableau} are analytic.
\end{lem}
\begin{proof}
    Follows immediately from the tableaux rules.
\end{proof}

\Cref{def:consistency} is auxiliary to seeing tableaux rules as a calculus for satisfiability. 

\begin{defi} \label{def:clash}\label{def:saturated}\label{def:consistency}
    A branch $\Theta$ of a tableau has a \emph{clash} iff one of the following conditions hold:
        \begin{enumerate}
            \item $\{i{:}\varphi,  i{:}\neg \varphi\}\subseteq\Theta$,
            \item $\{\tup{@_i =_\cmp @_j},\tup{@_i \neq_\cmp @_j}\}\subseteq\Theta$.
    \end{enumerate}
    We say that $\Theta$ is \emph{closed} if it contains a clash, otherwise it is \emph{open}.
    The tableau is \emph{closed} iff all its branches are closed, otherwise it is \emph{open}. 
    Finally, we say that $\Theta$ is \emph{saturated} iff it is fully expanded according to the rules of the tableau system.
\end{defi}

The notion of clash tells us that a contradiction has been found in a branch of the tableau.
At the same time, the notion of an open and saturated branch tells us that such a branch is free of contradictions (i.e., contradictions are impossible in the branch).  
To check whether  a node expression $\varphi$ is satisfiable, 
we begin with a single node tableaux having $i{:}\varphi$ as its root and successively apply the rules in~\Cref{rules:basic,rules:nominals,rules:internalization,rules:path}.
We answer the question negatively if at some stage the tableau closes, and answer the question positively if we find an open and saturated branch. 
The correctness of~\Cref{def:consistency} as a calculus for satisfiability follows from the results in~\Cref{sec:soundness,sec:completeness}.
Moreover, the results in~\Cref{sec:termination} tell us that this calculus yields a decision procedure for the satisfiability problem.
Even further, the results in~\Cref{sec:complexity} tells us that the computational complexity of this calculus is relatively low, insofar as it is in the same class as the satisfiability problem for Basic Modal Logic~\cite{handbook06}.